\section{Related Work}
\label{sec:related}

\subsection{Backdoor Attacks}
Backdoor poisoning was established by Gu \etal{} with BadNets, a small
checkerboard patch stamped onto a fraction of training images and relabelled to a
target class \cite{badnets}. Liu \etal{} showed that the same effect can be
obtained by perturbing weights rather than data \cite{trojannn}, and Chen \etal{}
gave an early data-poisoning formulation with a learned trigger
\cite{blendattack}. Later work made the trigger stealthier: clean-label attacks
that require no label change \cite{cleanlabel}, hidden-trigger attacks that
optimise the poisoned input to be far from the trigger at test time
\cite{hidden}, warping-based triggers \cite{wanet}, learnable and imperceptible
triggers \cite{lira}, input-aware triggers \cite{inputaware}, and physically
realisable triggers \cite{physical}. In federated learning, a single malicious
client can inject a backdoor that survives aggregation
\cite{bagdasaryan,bhagoji,dba}. Surveys by Li \etal{} \cite{survey} and Goldblum
\etal{} \cite{goldblum} organise this space and stress that defence evaluation
is fragmented.

Backdoor poisoning is a special case of data poisoning, which has a longer
history: support-vector poisoning \cite{biggio}, certified defences against
poisoning \cite{steinhardt_cert}, poisoning of regression learning
\cite{jagielski}, and attacks that break data sanitization at scale
\cite{kohpoison}, including gradient-matching poisons that transfer across
models \cite{witchesbrew}. The lesson of that line of work, that a defence tuned
to one poisoning strength can fail at another, motivates the sweep over recall
and precision here.

The design space has since widened well beyond dirty-label patches: clean-label
poisoning of feature extractors \cite{poisonfrogs}, subpopulation poisoning
\cite{subpop}, frequency-domain triggers \cite{freq}, reflection triggers
\cite{reflection}, dynamic triggers \cite{dynamic}, and latent backdoors
\cite{latent}; contrastive encoders \cite{carlinicl}; and instruction-tuned
language models \cite{poisonllm}. Data-sanitization defences can themselves be
evaded by stronger poisoning \cite{kohpoison}, and BackdoorBench consolidates
implementations and evaluation protocols across this space
\cite{backdoorbench}.

Our threat model is the simplest and most studied one: an attacker controls a
fraction of a static training set and stamps a fixed, visible trigger. We use two
instances, a $4\times4$ checkerboard and a blended pattern, precisely because
they are canonical. The contribution of this paper is not a new attack but the
measurement of what happens after one is detected imperfectly.

\subsection{Detection}
Detection methods score inputs, features, or model parameters. Neural Cleanse
optimises a minimal trigger per class and flags a class whose trigger is
anomalously small \cite{nc}. ABS stimulates neurons and observes whether a
single output class dominates \cite{abs}. STRIP perturbs inputs with
superimposed patterns and measures prediction entropy \cite{strip}. Activation
clustering projects per-class activations and separates them into two clusters
\cite{ac}, and spectral signatures analyses the top singular direction of
per-class representations \cite{ss}. Fine-pruning observes that backdoored
neurons are dormant on clean inputs and prunes them \cite{finepruning}.
Post-training scanners continue to broaden coverage; MM-BD, for example, scores
classes by a maximum-margin statistic and does not assume a patch-shaped
trigger \cite{mmbd}.

All of these methods trade missed detections against false positives, and their
operating points depend on poison rate, architecture, and attack. Detection is
therefore not a preprocessing detail that can be assumed away: it defines the
input to the removal stage. Recent defences continue to improve the detector
side, including variational formulations that jointly model trigger and
class-conditional structure \cite{seal}. Our study is complementary: we take
detector output as given and ask how much of its error the removal stage can
absorb. To keep the manipulation interpretable we parameterise that output by
recall and precision rather than by a detector-specific threshold, and we anchor
the parameterisation to two real detectors in Section~\ref{sec:cliff}.

\subsection{Removal and Unlearning}
The gold standard for removal is retraining on the retained data
\cite{badnets,finepruning}. Because full retraining is expensive, a large
literature approximates it. SISA shards the training set, trains a constituent
model per shard, and unlearns by retraining only the affected shards
\cite{sisa}. Amnesiac unlearns by ascending the loss on the forget set
\cite{amnesiac}, Eternal Sunshine approximates the retrained weights in a
quadratic Fisher sense \cite{eternal}, and certified removal provides a
differential-privacy-style guarantee for convex or strongly convex models
\cite{certified}. SCRUB trains a student to match the teacher on the retain set
while disagreeing on the forget set \cite{scrub}. Influence functions and
gradient tracing estimate which training points mattered and can be used to
down-weight them \cite{influence,tracin}, and example-forgetting scores have been
proposed as a way to choose which points a model should forget \cite{toneva}.
Suriyakumar and Wilson characterise
the limits of approximate-removal guarantees \cite{limits}. Surveys of machine
unlearning \cite{unlearnsurvey} and domain-specific benchmarks
\cite{tofu,wmdp} have grown around this literature, and what an unlearning claim
should certify is now a central question for language models
\cite{unlearnllm}.

A parallel family acts during or immediately after training rather than on a
forget set: anti-backdoor learning separates poison from clean data while
training \cite{abl}, decoupled training withholds the trigger association from
the feature extractor \cite{decouple}, adversarial neuron pruning removes
trigger-sensitive units \cite{anp}, and attention distillation fine-tunes the
backdoored model on clean data \cite{nad}. Trigger-reverse-engineering scanners
reconstruct a candidate trigger before removal \cite{tabor}. These methods
reduce the need for an exact forget set, but each still needs either clean data,
a clean reference model, or a reconstructed trigger; none is evaluated under a
forget set whose recall is known to be below one.

Closest to our recipe are methods that edit the association rather than delete
the data. Zhao and Wressnegger study the symmetry between learning a backdoor
and removing one \cite{twosides}. Repair-based pipelines relabel flagged samples
and fine-tune, optionally with an inoculation step that presents the trigger
with its correct label. These methods are attractive under imperfect detection
because they do not require the forget set to be exact, but their inoculation
step needs the trigger. We quantify both sides of that trade: the tolerance the
recipe buys, and the knowledge it consumes.

Two 2025 results shape our interpretation. Pawelczyk \etal{} show that
approximate unlearning methods, including methods designed for large language
models, fail to remove the effect of poisoned data even when they appear to
succeed on standard unlearning metrics \cite{pawelczyk}. Hu \etal{} show that
benign relearning restores supposedly unlearned behaviour \cite{jogging}, and
the same persistence under benign retraining has been observed in language
models trained to exhibit deceptive behaviour \cite{sleeper}. Our
durability study is a vision-domain analogue of the second result, and our
overall conclusion is consistent with the first: unlearning metrics that do not
probe the attack after removal can be misleading.

\subsection{Evaluation Under Imperfect Detection}
The gap we address is evaluative. Removal papers usually construct the forget
set from ground-truth poison labels, or from a detector threshold that is tuned
until recall is perfect. Detector papers report recall and precision but stop
before removal. Studies that connect the two usually vary a single threshold and
report a single operating point, which makes it hard to separate the effect of
missed poison from the effect of deleted clean data. Controlled manipulation of
both quantities, across method families and across domains, is what our grid
provides. The tolerance $\rstar$ is intended as a reusable summary: it answers
``how good must the detector be for this removal method to work?'' in the same
units for every method.

\subsection{Calibration and Risk Control}
Calibrate-then-Remove, one of the two refinements we test, borrows the logic of
distribution-free risk control. Conformal prediction converts a score into a
prediction set with a user-specified coverage guarantee under exchangeability
\cite{vovk,angelopoulos}, and conformal risk control extends the guarantee from
coverage to an arbitrary bounded loss \cite{crc}. The same machinery can in
principle choose a removal threshold that controls the risk of a live backdoor. The
guarantee depends on a proxy that is exchangeable with the test trigger. Our
experiment tests whether a synthesised trigger can serve as that proxy; the
answer, on both attacks, is no.

\subsection{Positioning}
This paper is a measurement study with one attempted mechanism and an explicit
negative result. It does not claim a new attack, a new detector, or a deployable
defence. Its contributions are the controlled forget-set construction, the
recall cliff and its transfer across architectures and attacks, the tolerance
metric $\rstar$, the demonstration that trigger synthesis breaks the only method
that tolerates imperfect recall, and the replication of the cliff in a radio
modulation task. The repair recipe we study is a strong baseline and an upper
bound; we label it as such wherever it appears.
